\documentclass[12pt, a4paper]{ctexart}

%==================== 导言区：宏包加载 ====================
\usepackage{amsmath, amssymb, amsthm} % 数学公式与符号
\usepackage{geometry}                 % 页面设置
\usepackage{xcolor}                   % 颜色支持
\usepackage{fancyhdr}                 % 页眉页脚
\usepackage{enumitem}                 % 列表定制
\usepackage{tcolorbox}                %以此美化文本框
\usepackage{titlesec}                 % 标题格式调整

%==================== 页面布局设置 ====================
\geometry{left=2.5cm, right=2.5cm, top=2.5cm, bottom=2.5cm}

%==================== 自定义视觉样式 ====================

% 1. 定义分隔线样式
\newcommand{\TopSeparator}{
    \par\noindent\rule{\textwidth}{1.5pt}\vspace{0.5em} % 粗实线 ━━━━
}
\newcommand{\AnalysisSeparator}{
    \par\noindent\rule{\textwidth}{0.5pt}\vspace{0.2em}\hrule\vspace{0.5em} % 双线效果 ════
}

% 2. 定义红色解析环境
%在这个环境内的所有文字和公式都会自动变成红色
\newenvironment{RedAnalysis}{
    \par
    \color{red} % 设置后续字体为红色
}{
    \par
    \color{black} % 结束时恢复黑色
}

% 3. 标题格式微调
\titleformat{\section}{\large\bfseries\heiti}{}{0em}{}[\vspace{-0.5em}]
\titleformat{\subsection}{\bfseries\kaishu}{}{0em}{}

\newcommand{\choicesFour}[4]{
    \begin{enumerate}[label=\Alph*., itemsep=0pt, parsep=0pt, topsep=5pt, labelsep=0.5em]
        \begin{minipage}{0.22\linewidth} \item #1 \end{minipage}
        \begin{minipage}{0.22\linewidth} \item #2 \end{minipage}
        \begin{minipage}{0.22\linewidth} \item #3 \end{minipage}
        \begin{minipage}{0.22\linewidth} \item #4 \end{minipage}
    \end{enumerate}
}
\newcommand{\choicesTwo}[4]{
    \begin{enumerate}[label=\Alph*., itemsep=0pt, parsep=0pt, topsep=5pt, labelsep=0.5em]
        \begin{minipage}{0.45\linewidth} \item #1 \end{minipage}
        \begin{minipage}{0.45\linewidth} \item #2 \end{minipage}
        \begin{minipage}{0.45\linewidth} \item #3 \end{minipage}
        \begin{minipage}{0.45\linewidth} \item #4 \end{minipage}
    \end{enumerate}
}
\newcommand{\choices}[4]{
    \begin{enumerate}[label=\Alph*., itemsep=0pt, parsep=0pt, topsep=5pt, labelsep=0.5em]
        \item #1
        \item #2
        \item #3
        \item #4
    \end{enumerate}
}

%==================== 文档开始 ====================
\begin{document}

% 标题信息
\title{\textbf{第九章 多元微积分（解析）}}
\author{数学习题讲义}
\date{\today}
\maketitle

% --- 题目 1 ---
\TopSeparator
\section*{题目1}

\textbf{【题目重现】} \\
设二元函数 $ z = e^{xy} + xy $ ，则 $ \frac{\partial z}{\partial x}|_{(1,2)} = $ \underline{\hspace{1cm}}。
\choicesFour{$ e^{2}+1 $}{$ 2e^{2}+2 $}{$ e+1 $}{$ 2e+1 $}

\AnalysisSeparator

\begin{RedAnalysis}

\subsection*{一、逐步解析}

\textbf{第一步：问题分析} \\
先求偏导数，再代入点 $(1,2)$ 的值。求 $\frac{\partial z}{\partial x}$ 时，把 $y$ 视为常数。

\textbf{详细步骤} \\
1. \textbf{求偏导}：
   \[ \frac{\partial z}{\partial x} = \frac{\partial}{\partial x}(e^{xy} + xy) = e^{xy} \cdot y + y = y(e^{xy} + 1) \]
2. \textbf{代入点值}：
   将 $x=1, y=2$ 代入：
   \[ \frac{\partial z}{\partial x}|_{(1,2)} = 2(e^{1 \cdot 2} + 1) = 2(e^2 + 1) = 2e^2 + 2 \]

\textbf{第四步：最终答案} \\
B

\subsection*{二、核心公式}
1. \textbf{偏导数定义}：求对 $x$ 偏导时，$y$ 为常数。

\end{RedAnalysis}

\vspace{2em}

% --- 题目 2 ---
\TopSeparator
\section*{题目2}

\textbf{【题目重现】} \\
若可微函数 $ z = f(x, y) $ 在点 $ (x_{0}, y_{0}) $ 取得极小值，下列各项说法正确的是 \underline{\hspace{1cm}}。
\choicesTwo
{$ f(x_0,y) $ 在 $ y=y_{0} $ 处的导数大于 0}
{$ f(x_0,y) $ 在 $ y=y_{0} $ 处的导数等于 0}
{$ f(x_0,y) $ 在 $ y=y_{0} $ 处的导数小于 0}
{$ f(x_0,y) $ 在 $ y=y_{0} $ 处的导数不存在}

\AnalysisSeparator

\begin{RedAnalysis}

\subsection*{一、逐步解析}

\textbf{第一步：问题分析} \\
极值存在的必要条件。若函数可微且在某点取得极值，则该点处的偏导数均为0。
$f(x_0, y)$ 可以看作关于 $y$ 的一元函数 $g(y)$。

\textbf{详细步骤} \\
1. 如果 $z=f(x,y)$ 在 $(x_0, y_0)$ 取得极小值，那么它在该点的两个偏导数都必须为0（如果是可微函数）。
2. 即 $f'_x(x_0, y_0) = 0$ 和 $f'_y(x_0, y_0) = 0$。
3. $f'_y(x_0, y_0)$ 的几何意义就是一元函数 $g(y) = f(x_0, y)$ 在 $y=y_0$ 处的导数。
4. 所以 $f(x_0, y)$ 在 $y=y_0$ 处的导数等于 0。

\textbf{第四步：最终答案} \\
B

\subsection*{二、核心公式}
1. \textbf{极值必要条件}：$\frac{\partial z}{\partial x} = 0, \frac{\partial z}{\partial y} = 0$。

\end{RedAnalysis}

\vspace{2em}

% --- 题目 3 ---
\TopSeparator
\section*{题目3}

\textbf{【题目重现】} \\
对于函数 $ f(x,y) $ 的每一个驻点 $ (x_{0},y_{0}) $ ，令 $ A=f_{xx}(x_{0},y_{0}) $ ， $ B=f_{xy}(x_{0},y_{0}) $ ， $ C = f_{yy}(x_{0}, y_{0}) $ , 若 $ AC - B^{2} < 0 $ , 则函数 \underline{\hspace{1cm}}。
\choicesFour{有极大值}{有极小值}{没有极值}{不定}

\AnalysisSeparator

\begin{RedAnalysis}

\subsection*{一、逐步解析}

\textbf{第一步：问题分析} \\
二元函数极值的充分条件判别法（Hessian矩阵行列式）。

\textbf{详细步骤} \\
1. 判别式 $\Delta = AC - B^2$。
2. 若 $\Delta > 0$，有极值（$A>0$极小，$A<0$极大）。
3. 若 $\Delta < 0$，无极值（该点为鞍点）。
4. 若 $\Delta = 0$，失效。

\textbf{第四步：最终答案} \\
C

\subsection*{二、核心公式}
1. \textbf{极值充分条件}：$AC-B^2<0 \Rightarrow$ 无极值。

\end{RedAnalysis}

\vspace{2em}

% --- 题目 4 ---
\TopSeparator
\section*{题目4}

\textbf{【题目重现】} \\
多元函数 $ f(x, y) $ 在点 $ (x_{0}, y_{0}) $ 处关于 $ y $ 的偏导数 $ f_{y}^{\prime}(x_{0}, y_{0}) = $ \underline{\hspace{1cm}}。
\choicesTwo
{$ \lim_{\Delta x \to 0} \frac{f(x_0 + \Delta x, y_0) - f(x_0, y_0)}{\Delta x} $}
{$ \lim_{\Delta x \to 0} \frac{f(x_0 + \Delta x, y_0 + \Delta y) - f(x_0, y_0)}{\Delta x} $}
{$ \lim_{\Delta y \to 0} \frac{f(x_0, y_0 + \Delta y) - f(x_0, y_0)}{\Delta y} $}
{$ \lim_{\Delta x \to 0} \frac{f(x_0 + \Delta x, y_0 + \Delta y) - f(x_0, y_0)}{\Delta y} $}

\AnalysisSeparator

\begin{RedAnalysis}

\subsection*{一、逐步解析}

\textbf{第一步：问题分析} \\
偏导数的定义。$f'_y$ 表示 $x$ 固定，$y$ 变化时的变化率。

\textbf{详细步骤} \\
1. 对 $y$ 的偏导数：固定 $x=x_0$，让 $y$ 从 $y_0$ 变到 $y_0+\Delta y$。
2. 定义式：
   \[ \lim_{\Delta y \to 0} \frac{f(x_0, y_0 + \Delta y) - f(x_0, y_0)}{\Delta y} \]
3. 观察选项：
   A. 是对 $x$ 的偏导。
   C. 是对 $y$ 的偏导。

\textbf{第四步：最终答案} \\
C

\subsection*{二、核心公式}
1. \textbf{定义}：$f_y(x_0,y_0) = \lim_{\Delta y \to 0} \frac{f(x_0, y_0+\Delta y)-f(x_0,y_0)}{\Delta y}$。

\end{RedAnalysis}

\vspace{2em}

% --- 题目 5 ---
\TopSeparator
\section*{题目5}

\textbf{【题目重现】} \\
若 $ f(x,y)=2x^{2}+y $ , 则 $ f_{x}^{'}= $ \underline{\hspace{1cm}}。
\choicesFour{4}{0}{2}{-1}

\AnalysisSeparator

\begin{RedAnalysis}

\subsection*{一、逐步解析}

\textbf{第一步：问题分析} \\
求偏导数 $f'_x$。

\textbf{详细步骤} \\
1. $f(x,y) = 2x^2 + y$。
2. $f_x(x,y) = 4x$。
3. $f_{xx}(x,y) = 4$。
   （注：若题目求 $f_x$，结果为 $4x$。若求 $f_{xx}$，结果为 4。若函数为 $f=2x+y$，则 $f_x=2$。
   根据选项判断，可能考察二阶偏导或一次函数。若为二阶偏导，选4。）

\textbf{第四步：最终答案} \\
$4x$ (若问二阶导则为4)

\subsection*{二、核心公式}
1. \textbf{求导法则}。

\end{RedAnalysis}

\vspace{2em}

% --- 题目 6 ---
\TopSeparator
\section*{题目6}

\textbf{【题目重现】} \\
求极限 $ \lim_{(x,y)\to(0,0)}\frac{\sqrt{xy+1}-1}{xy}= $ \underline{\hspace{2cm}}。

\AnalysisSeparator

\begin{RedAnalysis}

\subsection*{一、逐步解析}

\textbf{第一步：问题分析} \\
零比零型极限。分子有根式，进行有理化，或者利用等价无穷小。

\textbf{详细步骤} \\
1. 令 $u = xy$。当 $(x,y) \to (0,0)$ 时，$u \to 0$。
2. 原极限转化为一元极限：
   \[ \lim_{u \to 0} \frac{\sqrt{u+1}-1}{u} \]
3. 使用等价无穷小：$(1+u)^\alpha - 1 \sim \alpha u$。
   这里 $\alpha = 1/2$。
   $\sqrt{1+u} - 1 \sim \frac{1}{2}u$。
4. 代入：
   \[ \lim_{u \to 0} \frac{\frac{1}{2}u}{u} = \frac{1}{2} \]

\textbf{第四步：最终答案} \\
$ \frac{1}{2} $

\subsection*{二、核心公式}
1. \textbf{等价无穷小}：$\sqrt{1+x}-1 \sim \frac{1}{2}x$。

\end{RedAnalysis}

\vspace{2em}

% --- 题目 7 ---
\TopSeparator
\section*{题目7}

\textbf{【题目重现】} \\
函数 $ z = \sqrt{y - x^{2} + 1} $ 定义域为 \underline{\hspace{2cm}}。

\AnalysisSeparator

\begin{RedAnalysis}

\subsection*{一、逐步解析}

\textbf{第一步：问题分析} \\
根号下必须非负。

\textbf{详细步骤} \\
1. $y - x^2 + 1 \ge 0$。
2. 即 $y \ge x^2 - 1$。
   这是一个抛物线 $y=x^2-1$ 及其上方的区域。

\textbf{第四步：最终答案} \\
$ \{(x,y) | y \ge x^2 - 1\} $

\subsection*{三、考点分析}
\begin{itemize}
    \item \textbf{知识点归类}：函数定义域。
\end{itemize}

\end{RedAnalysis}

\vspace{2em}

% --- 题目 8 ---
\TopSeparator
\section*{题目8}

\textbf{【题目重现】} \\
设 $ z = x^{y} $ ，则 $ \frac{\partial z}{\partial x} = $ \underline{\hspace{2cm}}。

\AnalysisSeparator

\begin{RedAnalysis}

\subsection*{一、逐步解析}

\textbf{第一步：问题分析} \\
求对 $x$ 的偏导，将 $y$ 视为常数。此时是幂函数 $x^a$ 的形式。

\textbf{详细步骤} \\
1. $z = x^y$。
2. $\frac{\partial z}{\partial x} = y \cdot x^{y-1}$。

\textbf{第四步：最终答案} \\
$ yx^{y-1} $

\subsection*{二、核心公式}
1. \textbf{幂函数求导}：$(x^n)' = n x^{n-1}$。

\end{RedAnalysis}

\vspace{2em}

% --- 题目 9 ---
\TopSeparator
\section*{题目9}

\textbf{【题目重现】} \\
设 $ f(x, y, z) = x^{y}y^{z} $ ，则 $ \frac{\partial f}{\partial y} = $ \underline{\hspace{2cm}}。

\AnalysisSeparator

\begin{RedAnalysis}

\subsection*{一、逐步解析}

\textbf{第一步：问题分析} \\
对 $y$ 求偏导，$x, z$ 视为常数。
$f = x^y \cdot y^z$。两项都含有 $y$，使用乘法法则。
第一项 $x^y$ 是指数函数形式（底数常数）。
第二项 $y^z$ 是幂函数形式（指数常数）。

\textbf{详细步骤} \\
1. \[ \frac{\partial f}{\partial y} = \frac{\partial}{\partial y}(x^y) \cdot y^z + x^y \cdot \frac{\partial}{\partial y}(y^z) \]
2. $\frac{\partial}{\partial y}(x^y) = x^y \ln x$ （指数函数求导）。
3. $\frac{\partial}{\partial y}(y^z) = z y^{z-1}$ （幂函数求导）。
4. 合并：
   \[ = (x^y \ln x) y^z + x^y (z y^{z-1}) \]
   \[ = x^y y^z \ln x + x^y z y^{z-1} \]
   \[ = x^y y^{z-1} (y \ln x + z) \]

\textbf{第四步：最终答案} \\
$ x^y y^{z-1} (y \ln x + z) $

\subsection*{二、核心公式}
1. \textbf{乘法法则}：$(uv)' = u'v + uv'$。
2. \textbf{指数/幂导数}：$(a^x)' = a^x \ln a$; $(x^a)' = ax^{a-1}$。

\end{RedAnalysis}

\vspace{2em}

% --- 题目 10 ---
\TopSeparator
\section*{题目10}

\textbf{【题目重现】} \\
已知函数 $ z = x^{2} \arctan(2y) $ ，则全微分 $ dz = $ \underline{\hspace{2cm}}。

\AnalysisSeparator

\begin{RedAnalysis}

\subsection*{一、逐步解析}

\textbf{第一步：问题分析} \\
全微分公式 $dz = \frac{\partial z}{\partial x}dx + \frac{\partial z}{\partial y}dy$。

\textbf{详细步骤} \\
1. 求 $\frac{\partial z}{\partial x}$：
   \[ 2x \arctan(2y) \]
2. 求 $\frac{\partial z}{\partial y}$：
   \[ x^2 \cdot \frac{1}{1+(2y)^2} \cdot 2 = \frac{2x^2}{1+4y^2} \]
3. 组合：
   \[ dz = 2x \arctan(2y) dx + \frac{2x^2}{1+4y^2} dy \]

\textbf{第四步：最终答案} \\
$ 2x \arctan(2y) dx + \frac{2x^2}{1+4y^2} dy $

\subsection*{二、核心公式}
1. \textbf{全微分}：$dz = Z_x dx + Z_y dy$。

\end{RedAnalysis}

\vspace{2em}

% --- 题目 11 ---
\TopSeparator
\section*{题目11}

\textbf{【题目重现】} \\
设 $ u = e^{\frac{x}{y}} $ , 求 $ \frac{\partial^{2}u}{\partial x \partial y} $。

\AnalysisSeparator

\begin{RedAnalysis}

\subsection*{一、逐步解析}

\textbf{第一步：问题分析} \\
混合偏导数。先对 $x$ 求导，再对 $y$ 求导。

\textbf{详细步骤} \\
1. 对 $x$ 求导：
   \[ \frac{\partial u}{\partial x} = e^{\frac{x}{y}} \cdot \frac{1}{y} \]
2. 再对 $y$ 求导：
   \[ \frac{\partial}{\partial y} (\frac{1}{y} e^{\frac{x}{y}}) \]
   使用乘法法则：
   \[ = (\frac{1}{y})' e^{\frac{x}{y}} + \frac{1}{y} (e^{\frac{x}{y}})'_y \]
   \[ = (-\frac{1}{y^2}) e^{\frac{x}{y}} + \frac{1}{y} (e^{\frac{x}{y}} \cdot (-\frac{x}{y^2})) \]
   \[ = -\frac{1}{y^2} e^{\frac{x}{y}} - \frac{x}{y^3} e^{\frac{x}{y}} \]
   \[ = - \frac{e^{x/y}}{y^2} (1 + \frac{x}{y}) = - \frac{x+y}{y^3} e^{x/y} \]

\textbf{第四步：最终答案} \\
$ - \frac{x+y}{y^3} e^{x/y} $

\subsection*{二、核心公式}
1. \textbf{混合偏导数}：顺序不影响结果（连续时）。

\end{RedAnalysis}

\vspace{2em}

% --- 题目 12 ---
\TopSeparator
\section*{题目12}

\textbf{【题目重现】} \\
求由方程 $ e^{z}-xyz=0 $ 所确定的二元函数 $ z=f(x,y) $ 的全微分。

\AnalysisSeparator

\begin{RedAnalysis}

\subsection*{一、逐步解析}

\textbf{第一步：问题分析} \\
隐函数全微分。可以直接对方程两边微分。

\textbf{详细步骤} \\
1. 两边微分：
   \[ d(e^z) - d(xyz) = 0 \]
   \[ e^z dz - (yz dx + xz dy + xy dz) = 0 \]
2. 整理含有 $dz$ 的项：
   \[ (e^z - xy) dz = yz dx + xz dy \]
3. 解出 $dz$：
   \[ dz = \frac{yz}{e^z - xy} dx + \frac{xz}{e^z - xy} dy \]
   也可代入方程 $e^z = xyz$ 化简分母，但通常保留形式即可。

\textbf{第四步：最终答案} \\
$ dz = \frac{yz}{e^z - xy} dx + \frac{xz}{e^z - xy} dy $

\subsection*{二、核心公式}
1. \textbf{隐函数微分法}：两边直接微分后解出 $dz$。

\end{RedAnalysis}

\vspace{2em}

% --- 题目 13 ---
\TopSeparator
\section*{题目13}

\textbf{【题目重现】} \\
求函数 $ f(x,y)=\frac{1}{3}x^{3}+2x-3xy+\frac{3}{2}y^{2} $ 的极值，并判断是极大值还是极小值。

\AnalysisSeparator

\begin{RedAnalysis}

\subsection*{一、逐步解析}

\textbf{第一步：问题分析} \\
1. 求驻点（一阶偏导为0）。
2. 用二阶偏导判别法判断。

\textbf{详细步骤} \\
1. \textbf{求偏导}：
   $f_x = x^2 + 2 - 3y$
   $f_y = -3x + 3y$
2. \textbf{求驻点}：令 $f_x=0, f_y=0$。
   $-3x + 3y = 0 \Rightarrow y=x$。
   代入第一个方程：$x^2 + 2 - 3x = 0 \Rightarrow (x-1)(x-2) = 0$。
   解得 $x_1=1, x_2=2$。
   驻点：$M_1(1,1), M_2(2,2)$。
3. \textbf{求二阶偏导}：
   $f_{xx} = 2x, f_{xy} = -3, f_{yy} = 3$。
   $AC - B^2 = (2x)(3) - (-3)^2 = 6x - 9$。
4. \textbf{判别}：
   - 点 $M_1(1,1)$：$AC-B^2 = 6(1)-9 = -3 < 0$。
     无极值（鞍点）。
   - 点 $M_2(2,2)$：$AC-B^2 = 6(2)-9 = 3 > 0$。
     有极值。且 $A = 2(2) = 4 > 0$，为极小值。
5. \textbf{求极值}：
   $f(2,2) = \frac{1}{3}(8) + 4 - 12 + 6 = \frac{8}{3} - 2 = \frac{2}{3}$。

\textbf{第四步：最终答案} \\
极小值 $\frac{2}{3}$ （在点 $(2,2)$ 处）。

\subsection*{二、核心公式}
1. \textbf{极值判别}：$AC-B^2$。

\end{RedAnalysis}

\vspace{2em}

% --- 题目 14 ---
\TopSeparator
\section*{题目14}

\textbf{【题目重现】} \\
设 $ z = z(x, y) $ 是由 $ x^{2}z + 2y^{2}z^{2} + y = 0 $ 确定的函数，求 $ \frac{\partial z}{\partial y} $。

\AnalysisSeparator

\begin{RedAnalysis}

\subsection*{一、逐步解析}

\textbf{第一步：问题分析} \\
隐函数求导。方程两边对 $y$ 求导，$x$ 视为常数，$z$ 是 $y$ 的函数。

\textbf{详细步骤} \\
1. 对方程 $x^2 z + 2y^2 z^2 + y = 0$ 两边求偏导 $\frac{\partial}{\partial y}$：
   \[ x^2 \frac{\partial z}{\partial y} + 2(2y z^2 + y^2 \cdot 2z \frac{\partial z}{\partial y}) + 1 = 0 \]
2. 展开：
   \[ x^2 z_y + 4y z^2 + 4y^2 z z_y + 1 = 0 \]
3. 提取 $z_y$：
   \[ z_y (x^2 + 4y^2 z) = - (4y z^2 + 1) \]
4. 解出：
   \[ z_y = - \frac{4y z^2 + 1}{x^2 + 4y^2 z} \]

\textbf{第四步：最终答案} \\
$ - \frac{4y z^2 + 1}{x^2 + 4y^2 z} $

\subsection*{二、核心公式}
1. \textbf{隐函数求导公式}：$z_y = - \frac{F_y}{F_z}$。

\end{RedAnalysis}

\vspace{2em}

% --- 题目 15 ---
\TopSeparator
\section*{题目15}

\textbf{【题目重现】} \\
已知 $ z = f(xy, 2x + 3y) $ ，其中 $ f(u, v) $ 具有连续偏导数，求 $ \frac{\partial z}{\partial x} $。

\AnalysisSeparator

\begin{RedAnalysis}

\subsection*{一、逐步解析}

\textbf{第一步：问题分析} \\
复合函数求导（链式法则）。
设 $u = xy, v = 2x+3y$。

\textbf{详细步骤} \\
1. 公式：$\frac{\partial z}{\partial x} = \frac{\partial f}{\partial u} \frac{\partial u}{\partial x} + \frac{\partial f}{\partial v} \frac{\partial v}{\partial x}$。
2. 计算中间变量导数：
   $\frac{\partial u}{\partial x} = y$
   $\frac{\partial v}{\partial x} = 2$
3. 代入：
   $\frac{\partial z}{\partial x} = y f'_1 + 2 f'_2$ （或写为 $y f'_u + 2 f'_v$）。

\textbf{第四步：最终答案} \\
$ y f'_1 + 2 f'_2 $

\subsection*{二、核心公式}
1. \textbf{链式法则}：$z_x = f_u u_x + f_v v_x$。

\end{RedAnalysis}

\end{document}
